\begin{tabular}{llrrr}
\toprule
\textbf{Punto} & \textbf{Fuente} & \textbf{Intervalos} & \textbf{Dispersion [\%]} & \textbf{Sin el mas angosto [\%]} \\
\midrule
A @ 1600 V & Cs-137 & 5 de 6 & 1.0 & 1.0 \\
B @ 1600 V & Cs-137 & 5 de 6 & 49.5 & 1.4 \\
C @ 1600 V & Cs-137 & 5 de 6 & 1.6 & 1.6 \\
D @ 1600 V & Cs-137 & 4 de 6 & 0.8 & 0.5 \\
\midrule
\textit{paper, Fig. 5} & \textit{Cs-137} & 6 de 6 & 1.2 & 0.3 \\
\bottomrule
\end{tabular}
\par\smallskip
\begin{minipage}{0.95\linewidth}\footnotesize El mismo test de la Fig. 5 del paper: la misma corrida ajustada en seis intervalos anidados, definidos en unidades del canal del borde y tomados de sus propios paneles. Un intervalo con menos de 30 bines se saltea, igual que en la replica. La fila del paper sale de los seis valores que rotula su figura.\end{minipage}
